\setcounter{exercisegroup}{5}
\edef\CurrentExerciseGroup{\arabic{exercisegroup}}
\section*{\texorpdfstring{\S\CurrentExerciseGroup}{Section \CurrentExerciseGroup}}
\phantomsection
\addcontentsline{toc}{section}{\texorpdfstring{Exercises for \S\CurrentExerciseGroup}{Exercises for Section \CurrentExerciseGroup}}

\begin{enumerate}
\def\labelenumi{\arabic{enumi})}
\item
  Let \(\mathbf{T}\) and \(\mathbf{X}\) be two locally compact spaces, \(\mu\) a measure \(\geqslant 0\) on \(\mathbf{T}\) , \(\Lambda : t \mapsto \lambda_t\) a \(\mu\) -adequate mapping of \(\mathbf{T}\) into \(\mathcal{M}_+(X)\) , and \(g\) a function locally integrable for the measure \(\nu = \int \lambda_t d\mu(t)\) . Give an example showing that if \(\mathbf{X}\) is not countable at infinity, \(g\) may not be locally \(\lambda_t\) -integrable for any value of \(t\) (cf.~§3, Exer. 4).
\item
  Let \(g\) be a positive function on \(\mathbf{T}\) , measurable and essentially bounded for \(\mu\) , and let \(\nu = g \cdot \mu\) . Show that if a function \(\mathbf{f}\) defined on \(\mathbf{T}\) , with values in a Banach space or in \(\overline{\mathbf{R}}\) . is \(\mu\) -integrable, then \(\mathbf{f}\) is \(\nu\) -integrable. (Observe that \(\nu \leqslant a\mu\) for some constant \(a > 0\) , and make use of Prop. 15 of Ch. IV, §1, No.~3.)
\item
  Let X be a locally compact space, \(\Lambda : t \mapsto \lambda_t\) and \(\Lambda': t \mapsto \lambda_t'\) two \(\mu\) -adequate mappings of T into \(\mathcal{M}_+(X)\) , and let \(\nu = \int \lambda_t d\mu(t)\) , \(\nu' = \int \lambda_t' d\mu(t)\) . Show that if, locally almost everywhere for \(\mu\) , \(\lambda_t'\) is a measure with base \(\lambda_t\) , then \(\nu'\) is a measure with base \(\nu\) . (Consider a \(\nu\) -negligible compact subset K of X. It is \(\lambda_t\) -negligible locally almost everywhere, hence \(\lambda_t'\) -negligible locally almost everywhere. Apply Prop. 5 of §3.)
\item
  \begin{enumerate}
  \def\labelenumii{\alph{enumii})}
  \tightlist
  \item
    Let \(\mu\) be a moderated measure \(\geqslant 0\) on T. Show that there exists a function \(h \geqslant 0\) , continuous on T, such that \(h \cdot \mu\) is bounded and is equivalent to \(\mu\) (argue as for Prop. 11).
  \end{enumerate}
\end{enumerate}

\begin{enumerate}
\def\labelenumi{\alph{enumi})}
\setcounter{enumi}{1}
\tightlist
\item
  Let \((\mu_{n})\) be a sequence of moderated measures \(\geqslant0\) on T. Show that there exists a measure \(\mu\geqslant0\) on T such that each \(\mu_{n}\) has base \(\mu\) (reduce to the case that each of the \(\mu_{n}\) is bounded).
\end{enumerate}

\begin{enumerate}
\def\labelenumi{\arabic{enumi})}
\setcounter{enumi}{4}
\item
  Let \(\rho, \sigma\) be two atomic measures on T, and let M,N be the smallest sets carrying \(|\rho|\) and \(|\sigma|\) , respectively. For \(\rho\) and \(\sigma\) to be alien, it is necessary and sufficient that M \(\cap\) N = \(\varnothing\) . From this, deduce an example of an atomic measure \(\nu\) on the interval I = {[}0,1{]} of R, such that I is the support of \(\nu^{+}\) and of \(\nu^{-}\) .
\item
  \begin{enumerate}
  \def\labelenumii{\alph{enumii})}
  \tightlist
  \item
    Let \(\mu\) be a positive measure on T. If \(A \subset T\) is universally measurable and not locally \(\mu\) -negligible, show that there exists a positive measure \(\nu\) carried by A such that \(\nu \neq 0\) and \(\nu \leqslant \mu\) . (Take \(\nu = \varphi_{K} \cdot \mu\) , where K is a compact subset of A that is not \(\mu\) -negligible.)
  \end{enumerate}
\end{enumerate}

\begin{enumerate}
\def\labelenumi{\alph{enumi})}
\setcounter{enumi}{1}
\item
  Let \(\mathbf{M}\) be a universally measurable subset of \(\mathbf{T}\) . In order that a positive measure \(\lambda\) on \(\mathbf{T}\) be carried by \(\mathbf{M}\) , it is necessary and sufficient that \(\lambda\) be alien to every positive measure carried by \(\mathbb{C}\mathbf{M}\) (make use of \(a\) ).
\item
  Deduce from \(b\) ) that the measures \(\rho\) on \(T\) , such that \(|\rho|\) is carried by \(M\) , form a band in \(\mathcal{M}(T)\) . This band is vaguely closed in \(\mathcal{M}(T)\) if \(M\) is a closed subset of \(T\) (Ch. III, §2, Prop. 6).
\item
  Show that Lebesgue measure on I = {[}0,1{]} is the vague limit of a sequence of atomic measures carried by a fixed countable set A ⊂ I (cf.~Ch. III, §2, Th. 1 and Prop. 13). The band of measures carried by A is therefore not vaguely closed.
\end{enumerate}

¶ 7) a) Let \(\mu\) be a positive measure on T, and let A be a \(\mu\) -integrable set such that \(\mu(A) > 0\) . Show that if, for every \(\mu\) -integrable set \(B \subset A\) , either \(\mu(B) = 0\) or \(\mu(B) = \mu(A)\) , then there exists a point \(a \in A\) such that \(\mu(\{a\}) = \mu(A)\) . (Consider the intersection of the compact sets \(K \subset A\) such that \(\mu(K) = \mu(A)\) ; show that it is nonempty, has measure equal to \(\mu(A)\) , and reduces to a single point.)

\begin{enumerate}
\def\labelenumi{\alph{enumi})}
\setcounter{enumi}{1}
\tightlist
\item
  Suppose that \(\mu\) is a diffuse measure. For every \(\mu\) -integrable set A such that \(\mu(A) > 0\) , and every \(\varepsilon > 0\) , show that there exists a \(\mu\) -integrable subset B of A such that \(0 < \mu(B) \leqslant \varepsilon\) (observe, with the help of \(a\) ), that there exists a \(\mu\) -integrable subset C of A such that \(0 < \mu(C) \leqslant \frac{1}{2} \mu(A)\) ). From this, deduce that when X runs over the set of \(\mu\) -integrable subsets of A, the set of values of \(\mu(X)\) is the closed interval \([0, \mu(A)]\) (for every number \(\beta\) such that \(0 < \beta < \mu(A)\) , let \(\gamma\) be the supremum of the measures of the measurable subsets X of A such that \(\mu(X) \leqslant \beta\) ; first show that \(\gamma = \beta\) , using the preceding result, then prove that there exists an increasing sequence \((X_n)\) of measurable subsets of A such that \(\lim_{n \to \infty} \mu(X_n) = \beta\) ).
\end{enumerate}

¶ 8) a) Let \(\nu\) be a positive atomic measure on T, and let A be a \(\nu\) -integrable subset of T. Show that the set of values of \(\nu(X)\) , as X runs over the set of \(\nu\) -integrable subsets of A, is closed in R. (Let M be the smallest set carrying \(\nu\) . Assuming \(A \cap M\) infinite and arranging the points of \(A \cap M\) in a sequence \((a_n)\) , consider the mapping \(\varphi\) of the product space \(\{0,1\}^N\) into R defined by \(\varphi((\varepsilon_n)) = \sum_n \varepsilon_n \nu(\{a_n\})\) , and show

that it is continuous.)

\begin{enumerate}
\def\labelenumi{\alph{enumi})}
\setcounter{enumi}{1}
\item
  Deduce from a) and Exer. 7 b) that if \(\mu\) is any positive measure on T, and A is a \(\mu\) -integrable subset of T, then the set of values of \(\mu(X)\) , as X runs over the set of \(\mu\) -integrable subsets of A, is closed in R. Extend this result to the case that \(\mu\) is any real measure on T.
\item
  Deduce from Exer. 7 b) that if \(\mu\) is a diffuse real measure on T, then the set of values of \(\mu(X) = \mu^{+}(X) - \mu^{-}(X)\) , where X runs over the set of \(|\mu|\) -integrable subsets of T, is a closed interval of R (bounded or not).
\item
  Give an example of an atomic positive measure \(\nu\) on a noncompact locally compact space \(T\) , such that the set of values of \(\nu(X)\) , where \(X\) runs over the set of \(\nu\) -integrable subsets of \(T\) , is not closed (take \(\nu\) to be such that \(\inf_{t \in T} \nu(\{t\}) > 0\) ).
\end{enumerate}

\begin{enumerate}
\def\labelenumi{\arabic{enumi})}
\setcounter{enumi}{8}
\item
  Let \(\mu\) and \(\nu\) be two alien positive measures on T. Show that, for every number \(p\) such that \(1 \leqslant p \leqslant +\infty\) , the topological vector space \(\mathrm{L}^p(\mathrm{T}, \mu + \nu)\) is isomorphic to the product space of the topological vector spaces \(\mathrm{L}^p(\mathrm{T}, \mu)\) and \(\mathrm{L}^p(\mathrm{T}, \nu)\) .
\item
  \begin{enumerate}
  \def\labelenumii{\alph{enumii})}
  \tightlist
  \item
    Let \(\mu\) be a diffuse positive measure \(\neq 0\) on a locally compact space \(T\) . Show that there exists in \(\mathcal{L}^1(T, \mu)\) a sequence of functions \(f_n\) such that \(N_1(f_n) = 1\) for all \(n\) and such that the sequence of diffuse measures \(f_n \cdot \mu\) converges vaguely to a point measure \(\varepsilon_a\) . From this, deduce that \(L^1(T, \mu)\) (hence also \(L^\infty(T, \mu)\) ) is not reflexive.
  \end{enumerate}
\end{enumerate}

\begin{enumerate}
\def\labelenumi{\alph{enumi})}
\setcounter{enumi}{1}
\item
  Let \(\nu\) be an atomic positive measure on T, whose support is infinite. Let A be the smallest set carrying \(\nu\) , B a countable infinite subset of A, and \(\lambda\) the measure \(\varphi_{\mathrm{B}} \cdot \nu\) . Show that the dual of \(\mathrm{L}^{\infty}(\mathrm{T}, \lambda)\) is not separable, hence cannot be isomorphic to \(\mathrm{L}^{1}(\mathrm{T}, \lambda)\) (cf.~TVS, I, §2, Exer. 1).
\item
  Deduce from a) and b) that, for a positive measure \(\mu\) on a locally compact space T to be such that \(L^{1}(T,\mu)\) is reflexive, it is necessary and sufficient that the support of \(\mu\) be finite (make use of Exer. 9).
\end{enumerate}

¶ 11) Let S be a compact Stone space (Ch. II, §1, Exer. 13 f)); a positive measure \(\mu\) on S is said to be normal if, for every increasing directed family \((f_{\alpha})_{\alpha \in A}\) of continuous functions on S that is bounded above in \(\mathcal{C}(S)\) , and whose supremum in the fully lattice-ordered space \(\mathcal{C}(S)\) is \(f\) (which supremum is not necessarily equal to the upper envelope of the family \((f_{\alpha})\) ), one has \(\mu(f) = \sup_{\alpha \in A} \mu(f_{\alpha})\) . A real measure \(\lambda\) on S is said to be

normal if the positive measures \(\lambda^{+}\) and \(\lambda^{-}\) are normal.

\begin{enumerate}
\def\labelenumi{\alph{enumi})}
\item
  For a positive measure \(\mu\) on S to be normal, it is necessary and sufficient that every nowhere dense subset A be \(\mu\) -negligible (to see that the condition is necessary, consider the sets in S containing A that are open and closed; to see that the condition is sufficient, observe that the upper envelope and the supremum in \(\mathcal{C}(S)\) of an increasing directed family \((f_{\alpha})\) that is bounded above are equal on the complement of a meager set (cf.~Ch. II, §1, Exer. 13 f)). From this, deduce that the support of \(\mu\) is both open and closed.
\item
  Let \(\mu\) be a normal positive measure on S, \(f\) a \(\mu\) -measurable numerical function, and \(g\) the largest lower semi-continuous function on S that is \(\leqslant f\) . Show that \(f\) and \(g\) are equal almost everywhere for \(\mu\) (cf.~Ch. IV, §5, Exer. 16 b)). From this, deduce that
\end{enumerate}

for every \(\mu\) -measurable subset A of S, A - A and \(\overline{A}\) -A are \(\mu\) -negligible.

\begin{enumerate}
\def\labelenumi{\alph{enumi})}
\setcounter{enumi}{2}
\tightlist
\item
  Show that, in the Banach space \(\mathcal{M}(\mathrm{S})\) of measures on S, the set of normal measures is a closed linear subspace and is a band for the fully lattice-ordered space structure of \(\mathcal{M}(\mathrm{S})\) .
\end{enumerate}

¶ 12) A compact Stone space H is said to be hyperstonian if the union of the supports of the normal positive measures on H (Exer. 11) is dense.

\begin{enumerate}
\def\labelenumi{\alph{enumi})}
\item
  Let T be a locally compact space, \(\mu\) a positive measure on T. Show that there exists an isomorphism \(u \mapsto \theta_{u}\) of the Banach space \(\mathrm{L}^{\infty}(\mathrm{T}, \mu)\) onto the space \(\mathcal{C}(\mathrm{H})\) of continuous real-valued functions on a hyperstonian space H, such that \(\theta_{u}^{+} = \theta_{u^{+}}\) (cf.~Prop. 14 and Ch. II, §1, Exer. 13).
\item
  Let H be a hyperstonian compact space, and let \((\mu_{\iota})_{\iota\in I}\) be a family of normal positive measures on H such that the union of the supports of the \(\mu_{\iota}\) is dense in H. Show that, for a set to be nowhere dense in H, it is necessary and sufficient that it be \(\mu_{\iota}\) -negligible for all \(\iota\in I\) (observe that if A is \(\mu_{\iota}\) -negligible, then so is \(\overline{A}\) (Exer. 11 b)). Deduce from this that in a hyperstonian space, every meager set is nowhere dense.
\item
  The hypotheses being the same as in \(b\) ), show that if \(f\) is a numerical function \(\mu_{\iota}\) -measurable for all \(\iota \in I\) , there exists a continuous function \(g\) on \(H\) such that \(f\) and \(g\) coincide on the complement of a nowhere dense set (use \(b\) ) and Exer. 11 \(b\) ).
\item
  Let H be a hyperstonian space. Show that there exists a family \((\mathrm{G}_{\alpha})_{\alpha\in\mathrm{A}}\) of nonempty open and closed subsets of H, pairwise disjoint, whose union T is dense in H, and such that for every \(\alpha\in A\) , there exists a normal positive measure \(\mu_{\alpha}\) on \(G_{\alpha}\) with support \(G_{\alpha}\) (apply Zorn's lemma to the set of families of normal positive measures on H whose supports are mutually disjoint, and use Exer. 11 a)). Let \(\mu\) be the measure on T whose restriction to each \(G_{\alpha}\) is \(\mu_{\alpha}\) (Ch. III, §2, Prop. 1). Show that the mapping that, to every function \(f\in\mathcal{C}(\mathrm{H})\) , makes correspond the class in \(\mathrm{L}^{\infty}(\mathrm{T},\mu)\) of its restriction to T, is an isomorphism of the Banach space \(\mathcal{C}(\mathrm{H})\) onto \(\mathrm{L}^{\infty}(\mathrm{T},\mu)\) (use c) to show that the mapping is surjective).
\end{enumerate}

¶ 13) Let T be a locally compact space, \(\mu\) a positive measure on T, and \((f_n)\) a sequence of functions in \(\mathcal{L}^1(\mathrm{T},\mu)\) ; set \(\mu_n = f_n \cdot \mu\) and, for every \(\mu\) -measurable set A,

\[
\mu_ {n} (\mathrm{A}) = \mu_ {n} ^ {+} (\mathrm{A}) - \mu_ {n} ^ {-} (\mathrm{A}) = \int_ {\mathrm{A}} f _ {n} d \mu .
\]

\begin{enumerate}
\def\labelenumi{\alph{enumi})}
\item
  Show that if \(\mathbf{T}\) is not compact, it can happen that the sequence \((f_n)\) is not bounded in \(\mathcal{L}^1 (\mathrm{T},\mu)\) but the sequence of integrals \(\int f_n g d\mu\) is bounded for every \(g\in \mathcal{K}(\mathrm{T})\) .
\item
  Show that if the sequence \((\mu_n(A))\) is bounded for every subset \(A\) of \(T\) reduced to a point and for every open set \(A\) such that the measure induced by \(\mu\) on the boundary of \(A\) has finite support, then the sequence \((f_n)\) is bounded in \(\mathcal{L}^1\) , or, what amounts to the same, the sequence of norms \((\|\mu_n\|)\) is bounded. (Show first that every point \(t_0\) of \(T\) admits an open neighborhood \(U\) such that the sequence of numbers \(|\mu_n|(U)\) is bounded. For this, one argues by contradiction by proving that, in the contrary case, one could construct a strictly increasing sequence \((n_k)\) of integers, a decreasing sequence \((U_k)\) of neighborhoods of \(t_0\) , and a sequence \((W_k)\) of open sets quadrable for \(\mu\) (Ch. IV, §5, Exer. 17) having the following properties: \(\overline{U}_k \subset U_{k-1}\) , \(\mu(U_k - \{t_0\}) \leqslant 1/k\) , \(|\mu_{n_i}|(U_k - \{t_0\}) \leqslant 1\) for \(i < k\) , \(\overline{W}_k \subset U_k - \overline{U}_{k+1}\) , and finally
\end{enumerate}

\[
\left| \mu_ {n _ {k}} \left(\mathrm{W} _ {k}\right) \right| > k + \sum_ {i <   k} \left| \mu_ {n _ {k}} \left(\mathrm{W} _ {i}\right) \right|.
\]

Finally, consider the union W of the \(W_{k}\) , to obtain a contradiction (the `gliding hump method'). Then show by an analogous argument that there exists a compact subset K of T such that the sequence \(|\mu_{n}|(T - K)\) is bounded.)

\begin{enumerate}
\def\labelenumi{\alph{enumi})}
\setcounter{enumi}{2}
\item
  Deduce from \(b\) ) that if, for every open set \(A\) of \(T\) , the sequence \((\mu_n(A))\) is bounded, then the sequence \((f_n)\) is bounded in \(\mathcal{L}^1\) .
\item
  On the interval \([0,1]\) , if one takes for \(\mu_{n}\) the measure defined by the mass n at the point t=0 and the mass -n at the point \(t=1/n\) , then the sequence \((\mu_{n}(A))\) is bounded for every open set A that is quadrable for all the measures \(|\mu_{n}|\) ; however, the sequence of norms \(\|\mu_{n}\|\) is not bounded. Similarly, if one takes for \(\mu_{n}\) the measure defined by the mass n at the point \(t=1/n\) and the mass -n at the point \(t=1/(n+1)\) , then the sequence \((\mu_{n}(A))\) is bounded for every interval A contained in \([0,1]\) (hence for every finite subset A) without the sequence \((\|\mu_{n}\|)\) being bounded.
\end{enumerate}

\begin{enumerate}
\def\labelenumi{\arabic{enumi})}
\setcounter{enumi}{13}
\tightlist
\item
  Let \(\theta\) be a linear form on the space \(\mathcal{L}^{\infty}(\mathrm{T},\mu)\) ; in order that \(\theta\) be of the type \(f\mapsto\int fg d\mu\) , where \(g\in\mathcal{L}^{1}(\mathrm{T},\mu)\) , it is necessary and sufficient that \(\theta\) satisfy the following conditions: \(1^{\circ}\) for every \(\varepsilon>0\) , there exists a \(\delta>0\) such that the relation \(\mu^{*}(\mathrm{A})\leqslant\delta\) implies \(|\theta(h)|\leqslant\varepsilon\) for every \(\mu\) -measurable function \(h\) such that \(|h|\leqslant\varphi_{\mathrm{A}}\) ; \(2^{\circ}\) for every \(\varepsilon>0\) , there exists a compact set K such that, for every \(\mu\) -integrable set B \(\subset\) T−K, one has \(|\theta(\varphi_{\mathrm{B}})|\leqslant\varepsilon\) . (Use the Lebesgue--Nikodym theorem.)
\end{enumerate}

¶ 15) Let H be a bounded subset of the space \(\mathcal{L}^1 (\mathrm{T},\mu)\) , \(\widetilde{\mathrm{H}}\) its canonical image in the Banach space \(\mathrm{L}^1 (\mathrm{T},\mu)\) . Show that the following conditions are equivalent:

\(\alpha)\) The following set of two conditions holds:

\(\alpha_{1})\) for every \(\varepsilon > 0\) , there exists a compact set \(L \subset T\) such that \(\int_{T - L} |f| d\mu \leqslant \varepsilon\) for all \(f \in H\) ;

\(\alpha_{2})\) for every \(\varepsilon > 0\) , there exists an \(\eta > 0\) such that, for every open set U of T such that \(\mu^{*}(\mathrm{U}) \leqslant \eta\) , one has \(\int_{\mathrm{U}} |f| d\mu \leqslant \varepsilon\) for all \(f \in \mathrm{H}\) .

\(\beta)\) Both condition \(\alpha_{1}\) ) and the following condition \(\beta_{2}\) ) hold:

\(\beta_{2})\) for every compact set \(\mathbf{K} \subset \mathbf{T}\) and every \(\varepsilon > 0\) , there exists an open neighborhood \(\mathbf{U}\) of \(\mathbf{K}\) such that \(\int_{\mathbf{U} - \mathbf{K}} |f| d\mu \leqslant \varepsilon\) for all \(f \in \mathbf{H}\) .

\(\gamma)\) For every sequence \((g_n)\) of functions in \(\mathcal{L}^{\infty}\) , uniformly bounded and convergent in measure to a function \(g\) , one has \(\lim_{n\to \infty}\int fg_n d\mu = \int fg d\mu\) uniformly as \(f\) runs over H.

\(\delta)\) For every sequence \((h_n)\) of functions continuous on T and tending to 0 at the point at infinity, uniformly bounded and such that \(\lim_{n\to \infty}h_n(t) = 0\) for every \(t\in \mathbb{T}\) , one has \(\lim_{n\to \infty}\int fh_n d\mu = 0\) uniformly as \(f\) runs over H.

\(\zeta)\) For every infinite sequence \((\mathbf{U}_n)\) of pairwise disjoint open sets of T, one has \(\lim_{n\to \infty}\int_{\mathbf{U}_n}f d\mu = 0\) uniformly as \(f\) runs over H.

(To show that \(\zeta\) ) implies \(\beta\) ) and that \(\beta\) ) implies \(\alpha\) ), use a `gliding hump method' in the same way as in Exer. 13 \(b\) ). Then prove that \(\alpha\) ) implies \(\gamma\) ), that \(\gamma\) implies \(\delta\) ), and that \(\delta\) ) implies \(\zeta\) .)

If it is assumed moreover that, for every \(t \in \mathbf{T}\) such that \(\mu(\{t\}) \neq 0\) , the set of \(f(t)\) is bounded in \(\mathbf{R}\) as \(f\) runs over \(\mathbf{H}\) , show that the preceding conditions are equivalent to the condition:

\(\theta)\) H is a relatively compact subset of \(\mathbf{L}^1\) for the weakened topology \(\sigma (\mathbf{L}^{1},\mathbf{L}^{\infty})\)

(To show that \(\theta\) ) implies \(\beta\) ), argue by contradiction, using Šmulian's theorem (TVS, IV, §5, No.~3, Th. 2) and a `gliding hump method'. To show that \(\alpha\) ) implies \(\theta\) ), show that H is bounded in \(\mathcal{L}^1\) , then apply Eberlein's theorem (TVS, IV, §5, No.~3, Th. 1), and make use of Exer. 14.)

¶ 16) a) Let \((f_n)\) be a sequence of functions in \(\mathcal{L}^1(\mathrm{T},\mu)\) . Show that the following conditions are equivalent:

\(\alpha)\) The sequence \((f_n)\) is convergent in \(\mathbf{L}^1\) for the weakened topology \(\sigma(\mathbf{L}^1, \mathbf{L}^\infty)\)

\(\beta)\) The set of \(f_{n}\) is relatively compact in \(\mathbf{L}^{1}\) for the weakened topology, and the sequence of measures \(f_{n} \cdot \mu\) converges vaguely in \(\mathcal{M}(\mathrm{T})\) .

\(\gamma)\) For every open subset U of T, the sequence of numbers \(\int_{\mathrm{U}} f_n d\mu\) is convergent in \(\mathbf{R}\) .

(To prove that \(\beta\) ) implies \(\alpha\) ), use criterion \(\alpha\) ) of Exer. 15 and the definition of measurable function. To prove that \(\gamma\) ) implies \(\beta\) ), consider first the special case of the space \(\mathrm{L}^1(\mathbf{N})\) (for the discrete measure defined by the mass \(+1\) at each point) using Exer. 15 of TVS, IV, §2. In the general case, apply the criterion \(\zeta\) ) of Exer. 15 to reduce to the case of \(\mathrm{L}^1(\mathbf{N})\) , by associating to each \(f_n\) the summable sequence ( \(\int_{\mathrm{U}_m} f_n d\mu)_{m \in \mathbf{N}}\) .)

\begin{enumerate}
\def\labelenumi{\alph{enumi})}
\setcounter{enumi}{1}
\item
  Deduce from these results that, in \(L^{1}\) , every Cauchy sequence for the weakened topology is convergent for that topology.
\item
  On the interval T = {[}0,1{]}, let \(\mu\) be Lebesgue measure and, for every integer \(n \geqslant 1\) , let \(\mu_{n}\) be the measure defined by the mass 1/n at each of the points k/n ( \(0 \leqslant k < n\) ); let \(\nu\) be a positive measure on T such that \(\mu = g \cdot \nu\) , \(\mu_{n} = f_{n} \cdot \nu\) (Exer. 4 b)). Show that the sequence \((\mu_{n})\) converges vaguely to \(\mu\) , \(\mu_{n}(A)\) tends to \(\mu(A)\) for every finite subset A of T and for every open set A such that the measure induced by \(\nu\) on the boundary of A has finite support, but that \(\mu_{n}(U)\) does not tend to \(\mu(U)\) for all of the open sets \(U \subset T\) (cf.~Exer. 13 b)).
\end{enumerate}

\begin{enumerate}
\def\labelenumi{\arabic{enumi})}
\setcounter{enumi}{16}
\tightlist
\item
  \begin{enumerate}
  \def\labelenumii{\alph{enumii})}
  \tightlist
  \item
    Let \(\mu\) be Lebesgue measure on the interval \(T = [0, +\infty[\) . Show that if \(1 \leqslant r < s < +\infty\) , the topologies induced on \(L^r \cap L^s\) by the weakened topologies of \(L^r\) and \(L^s\) are not comparable (cf.~Ch. IV, §6, Exer. 8). If \(f_n\) is the characteristic function of the interval \([n, n+1]\) , show that the sequence \((\widetilde{f}_n)\) tends to 0 for the weakened topology of all the \(L to the p\) such that \(p > 1\) , but not for the weakened topology of \(L^1\) .
  \end{enumerate}
\end{enumerate}

\begin{enumerate}
\def\labelenumi{\alph{enumi})}
\setcounter{enumi}{1}
\item
  Let \(\lambda\) be Lebesgue measure on the interval \(T = [0,1]\) . Show that if \(r < s\) , the weakened topology of \(L^s\) is strictly finer than the topology on \(L^s\) induced by the weakened topology of \(L^r\) (cf.~Ch. IV, §6, Exer. 8).
\item
  Let \(\mu\) be the positive measure on a discrete space for which every point has measure 1. Show that if \(1 \leqslant r < s < +\infty\) , the weakened topology of \(\mathbf{L}^r\) is strictly finer than the topology induced on \(\mathbf{L}^r\) by the weakened topology of \(\mathbf{L}^s\) (Ch. IV, §6, Exer. 9).
\end{enumerate}

\begin{enumerate}
\def\labelenumi{\arabic{enumi})}
\setcounter{enumi}{17}
\tightlist
\item
  Let A be a set contained in \(L^{r} \cap L^{s}\) that is bounded in both \(L^{r}\) and \(L^{s}\) ( \(1 < r < s < +\infty\) ). Show that for every p such that \(r \leqslant p \leqslant s\) , A is bounded in \(L^{p}\) , and that the topologies induced on A by the weakened topologies of the \(L^{p}\) are identical (note that \(\widetilde{\mathcal{K}}(\mathrm{T})\) is dense in all of the spaces \(L^{q}\) where \(1 \leqslant q < +\infty\) ). Show that the property ceases to be valid if r = 1 (cf.~Exer. 17 a)).
\end{enumerate}

¶ 19) Let \(\mu\) be a positive measure on T.

\begin{enumerate}
\def\labelenumi{\alph{enumi})}
\item
  Let \((f_n)\) be a sequence of functions in \(\mathcal{L}^p\) ( \(1 < p < +\infty\) ) bounded in \(\mathcal{L}^p\) and convergent in measure to 0. Show that the sequence \((\widetilde{f}_n)\) converges to 0 in \(\mathbf{L}^p\) for the weakened topology.
\item
  Give an example of a sequence \((f_n)\) of functions in \(\mathcal{L}^1\) , bounded in \(\mathcal{L}^1\) and convergent in measure to 0, but such that the sequence \((\widetilde{f}_n)\) does not converge to 0 for the weakened topology of \(\mathrm{L}^1\) (take \(\mu\) to be Lebesgue measure on \(\mathrm{T} = [0,1]\) , and use criterion \(\alpha\) ) of Exer. 15).
\item
  Let \(\mu\) be Lebesgue measure on \(T = [0,1]\) . Show that, if \(f_{n}(t) = \sin nt\) , the sequence \((\widetilde{f}_n)\) converges to 0 for the weakened topology of all the \(L^{p}\) such that \(1 \leqslant p < +\infty\) , but \(f_{n}\) does not converge in measure to 0.
\end{enumerate}

\begin{enumerate}
\def\labelenumi{\arabic{enumi})}
\setcounter{enumi}{19}
\tightlist
\item
  \begin{enumerate}
  \def\labelenumii{\alph{enumii})}
  \tightlist
  \item
    Let \((f_n)\) be a sequence of functions in \(\mathcal{L}^1\) such that: \(1^\circ\) the sequence \((f_n)\) converges in measure to \(0; 2^\circ\) for every \(\varepsilon > 0\) , there exists an \(\eta > 0\) such that the relation \(\mu^*(A) \leqslant \eta\) implies \(\limsup_{n \to \infty} \int_A^* |f_n| d\mu \leqslant \varepsilon; 3^\circ\) for every \(\varepsilon > 0\) , there exists an integrable set \(B\) such that \(\int_{T - B} |f_n| d\mu \leqslant \varepsilon\) for all \(n\) . Show that the sequence \((f_n)\) converges in mean to \(0\) .
  \end{enumerate}
\end{enumerate}

\begin{enumerate}
\def\labelenumi{\alph{enumi})}
\setcounter{enumi}{1}
\tightlist
\item
  Let \(\mu\) be Lebesgue measure on \(T = [0,1]\) . For \(1 < p < +\infty\) , denote by \(f_n\) the function equal to \(n^{2/p}\) in the interval \(\left[\frac{1}{n+1}, \frac{1}{n}\right]\) , and to 0 elsewhere. Show that the sequence \((f_n)\) satisfies the three conditions of \(a)\) and that the sequence \((\widetilde{f}_n)\) converges to 0 for the weakened topology of \(L to the p\) , but not for the topology of convergence in mean of order \(p\) .
\end{enumerate}

¶ 21) a) Let \((f_n)\) be a sequence of functions in \(\mathcal{L}^1\) , such that: \(1^\circ \liminf_{n \to \infty} f_n(t) \geqslant 0\) almost everywhere in T; \(2^\circ\) for every \(\varepsilon > 0\) , there exists an \(\eta > 0\) such that the relation \(\mu^*(A) \leqslant \eta\) implies \(\limsup_{n \to \infty} \int_A^* |f_n| d\mu \leqslant \varepsilon\) ; \(3^\circ\) for every \(\varepsilon > 0\) , there exists an integrable set B such that \(\int_{T-B} |f_n| d\mu \leqslant \varepsilon\) for all \(n\) ; \(4^\circ \lim_{n \to \infty} \int f_n d\mu = 0\) . Show that the sequence \((f_n)\) converges in mean to 0. (Observe that, for every \(\varepsilon > 0\) , if \(C_m\) is the set of \(t \in T\) such that \(f_n(t) \geqslant -\varepsilon\) for at least one \(n \geqslant m\) , then the sequence \((C_m)\) is decreasing and has negligible intersection.)

\begin{enumerate}
\def\labelenumi{\alph{enumi})}
\setcounter{enumi}{1}
\tightlist
\item
  Let \((f_n)\) be a sequence of functions in \(\mathcal{L}^1\) such that: \(1^\circ\) there exists an integrable function \(g\) such that \(f_n \geqslant g\) for all \(n\) ; \(2^\circ \liminf_{n \to \infty} f_n(t) \geqslant 0\) almost everywhere; \(3^\circ \limsup_{n \to \infty} \int f_n d\mu \leqslant 0\) . Show that the sequence \((f_n)\) converges in mean to 0. (First show that, for every measurable set A, \(\liminf_{n \to \infty} \int_{A} f_n d\mu \geqslant 0\) ; then observe that, for every measurable set A,
\end{enumerate}

\[
\limsup _ {n \to \infty} \int f _ {n} d \mu \geqslant \limsup _ {n \to \infty} \int_ {\mathrm{A}} f _ {n} d \mu + \liminf _ {n \to \infty} \int_ {\mathrm{T-A}} f _ {n} d \mu .
\]

¶ 22) In the Banach space \(\mathrm{E} = \mathcal{M}^1(\mathrm{T})\) of bounded measures on a locally compact space \(\mathrm{T}\) , let \(\mathrm{H}\) be a compact set for the weakened topology.

\begin{enumerate}
\def\labelenumi{\alph{enumi})}
\item
  Show that, for every \(\varepsilon > 0\) , there exists a positive measure \(\mu_{\varepsilon}\) on T such that every measure \(\nu \in H\) is the sum of a positive measure with base \(\mu_{\varepsilon}\) and a measure \(\lambda\) alien to \(\mu_{\varepsilon}\) and of norm \(\leqslant \varepsilon\) . (Make use of the fact that the space E, equipped with its norm and its order structure, is isomorphic to a space \(L^{1}(S, \rho)\) , where S is the union of compact Stone spaces and \(\rho\) is a positive measure on S (Ch. IV, §4, Exer. 10), and apply in \(L^{1}(S, \rho)\) the criterion \(\alpha_{1}\) ) of Exer. 15.)
\item
  Deduce from a) that there exists a positive measure \(\mu\) on T such that all of the measures \(\nu \in H\) have base \(\mu\) (use Exer. 4 b)).
\end{enumerate}

\begin{enumerate}
\def\labelenumi{\arabic{enumi})}
\setcounter{enumi}{22}
\item
  Let \(\mu\) be a positive measure on a locally compact space \(T\) , \(p\) a number such that \(1 \leqslant p < +\infty\) , and \(q\) the conjugate exponent. Show that if \(g\) is a finite measurable function such that, for every function \(f \in \mathcal{L}^p\) , the function \(fg\) is integrable, then \(g\) is locally almost everywhere equal to a function in \(\mathcal{L}^q\) . (Show that the mapping \(f \mapsto fg\) of \(\mathcal{L}^p\) into \(\mathcal{L}^1\) is continuous, using the closed graph theorem (TVS, I, §3, Cor. 5 of Th. 1).)
\item
  Let \(p, q\) be two conjugate exponents. If B (resp. C) is a bounded subset of \(L to the p\) (resp. \(L^q\) ), show that the mapping of \(B \times C\) into \(L^1\) , deduced from \((f, g) \mapsto fg\) by passage to quotients, is not necessarily continuous when \(L to the p\) , \(L^q\) and \(L^1\) are equipped with the topologies \(\sigma(L to the p, L^q)\) , \(\sigma(L^q, L to the p)\) and \(\sigma(L^1, L^\infty)\) , respectively (cf.~Exer. 10).
\end{enumerate}

¶ 25) a) If u and v are any two real numbers, prove, for \(1 < p < +\infty\) , the inequality

\[
| u + v | ^ {p} \leqslant | u | ^ {p} + p v u | u | ^ {p - 2} + a \sum_ {r = 2} ^ {[ p ]} | v | ^ {r} | u | ^ {p - r} + b | v | ^ {p},
\]

where \(a\) and \(b\) are two constants depending only on \(p\) , and \([p]\) is the integral part of \(p\) (FRV, III, §2, Exer. 6).

\begin{enumerate}
\def\labelenumi{\alph{enumi})}
\setcounter{enumi}{1}
\tightlist
\item
  Let \((f_n)\) be a sequence that converges to 0 in \(\mathbf{L}^p\) for the weakened topology \((1 < p < +\infty)\) . Show that there exists a subsequence \((f_{n_k})\) of \((f_n)\) such that, on setting \(s_m = \sum_{k=1}^{m} f_{n_k}\) , one has
\end{enumerate}

\[
\left| \int s _ {m - 1} | s _ {m - 1} | ^ {p - 2} f _ {n _ {m}} d \mu \right| \leqslant 1
\]

(define the sequence \((n_k)\) recursively). Deduce from this, using \(a\) and Hölder's inequality, that if \(p > 2\) there exist two constants \(a, b\) such that

\[
\mathrm{N} _ {1} (| s _ {n} | ^ {p}) \leqslant \mathrm{N} _ {1} (| s _ {n - 1} | ^ {p}) + a + b \mathrm{N} _ {1} (| s _ {n - 1} | ^ {p - 2}),
\]

and that if \(1 < p \leqslant 2\) there exists a constant \(c\) such that

\[
\mathrm{N} _ {1} \big (| s _ {n} | ^ {p} \big) \leqslant \mathrm{N} _ {1} \big (| s _ {n - 1} | ^ {p} \big) + c.
\]

Conclude that

\[
\begin{array}{l l} \mathrm{N} _ {p} (s _ {n}) = O (n ^ {1 / 2}) & \text { for } p > 2 \\ \mathrm{N} _ {p} (s _ {n}) = O (n ^ {1 / p}) & \text { for } 1 <   p \leqslant 2. \end{array}
\]

\begin{enumerate}
\def\labelenumi{\alph{enumi})}
\setcounter{enumi}{2}
\tightlist
\item
  Show that the preceding results cannot in general be improved upon (cf.~Exers. 19 and 20 b)).
\end{enumerate}

¶ 26) Let \(\mu_1, \ldots, \mu_m\) be a finite number of measures on a locally compact space \(T\) , which one can write in the form \(\mu_k = f_k \cdot \mu\) , where \(\mu\) is a measure \(\geqslant 0\) and \(|f_k| \leqslant 1\) (No.~9).

\begin{enumerate}
\def\labelenumi{\alph{enumi})}
\item
  Let \(\Phi\) be a set of functions defined and positively homogeneous on \(\mathbf{R}^m\) and such that \(u(f_1, \ldots, f_m)\) is locally \(\mu\) -integrable for every function \(u \in \Phi\) . Show that if a sequence \((u_n)\) of functions in \(\Phi\) converges uniformly on every compact subset of \(\mathbf{R}^m\) to a function \(u\) , then \(u(\mu_1, \ldots, \mu_m)\) is the limit of the sequence of measures \(u_n(\mu_1, \ldots, \mu_m)\) for the quasi-strong topology on \(\mathcal{M}(\mathrm{T})\) (Ch. III, §1, Exer. 8).
\item
  Suppose that \(u \in \Phi\) is continuous on \(\mathbf{R}^m\) . Show that the mapping
\end{enumerate}

\[
(\mu_ {1}, \dots , \mu_ {m}) \mapsto u (\mu_ {1}, \dots , \mu_ {m})
\]

of \((\mathcal{M}(\mathrm{T}))^m\) into \(\mathcal{M}(\mathrm{T})\) is continuous for the quasi-strong topology (begin by considering the case that \(u\) is Lipschitz; then use \(a\) ) and the Stone-Weierstrass theorem).

\begin{enumerate}
\def\labelenumi{\alph{enumi})}
\setcounter{enumi}{2}
\tightlist
\item
  Suppose that \(u\) is continuous on \(\mathbf{R}^m\) . Let \(g\) be any function in \(\mathcal{H}(\mathrm{T})\) , and \(K\) its support. For every \(\varepsilon > 0\) , show that there exists a finite open covering \((A_i)\) of \(K\) formed of relatively compact sets, such that for every finite open covering \((B_j)\) of \(K\) , finer than \((A_i)\) , and for every family \((h_j)\) of continuous mappings of \(T\) into {[}0,1{]} such that \(h_j\) has its support in \(B_j\) and such that \(\sum_{j} h_j(t) = 1\) on \(K\) , one has
\end{enumerate}

\[
\left| \int g \cdot u (f _ {1}, \dots , f _ {m}) d \mu - \sum_ {j} u \left(\int g h _ {j} f _ {1} d \mu , \dots , \int g h _ {j} f _ {m} d \mu\right) \right| \leqslant \varepsilon
\]

(consider first the case that \(u\) is Lipschitz and the \(f_k\) are continuous, then pass to the case that \(u\) is Lipschitz and the \(f_k\) are locally integrable, and finally the general case using the Stone-Weierstrass theorem).

\begin{enumerate}
\def\labelenumi{\alph{enumi})}
\setcounter{enumi}{3}
\tightlist
\item
  In \(\mathbf{R}^2\) , let \(u(x_1, x_2)\) be the positively homogeneous function equal to \(\sqrt{x_1^2 + x_2^2}\) when \(x_1 / x_2\) is irrational, and to 0 in the contrary case. Let \(\mu\) be Lebesgue measure on \(\mathrm{T} = [0, 1]\) , and let \(\nu\) be the measure on \(\mathrm{T}\) defined by \(d\nu(t) = t d\mu(t)\) . Show that if \(g\) is a continuous function \(\geqslant 0\) on \(\mathrm{T}\) , then \(\sum_{j} u\left( \int gh_j d\mu, \int gh_j d\nu \right)\) does not tend to any
\end{enumerate}

limit with respect to the directed set of finite open coverings of the support of g.

\begin{enumerate}
\def\labelenumi{\arabic{enumi})}
\setcounter{enumi}{26}
\tightlist
\item
  Let \(\mathbf{X}\) be a locally compact space, \(\Lambda : t \mapsto \lambda_t\) a scalarly essentially \(\mu\) -integrable mapping of \(\mathbf{T}\) into \(\mathcal{M}_+(\mathbf{X})\) . Suppose that \(\Lambda\) is pre-adequate for every measure \(\varphi_{\mathbf{K}} \cdot \mu\) , where \(\mathbf{K}\) is compact.
\end{enumerate}

\begin{enumerate}
\def\labelenumi{\alph{enumi})}
\item
  Show that \(\Lambda\) is pre-adequate for every measure \(\varphi_{\mathrm{A}}\cdot \mu\) , where \(\mathbf{A}\) is a \(\mu\) -measurable subset of \(\mathbf{T}\) (consider the compact subsets of \(\mathbf{A}\) ).
\item
  Show that \(\Lambda\) is pre-adequate for every measure \(f \cdot \mu\) , where \(f\) is positive, \(\mu\) -measurable, and takes on only finitely many values.
\item
  Show that \(\Lambda\) is pre-adequate for every measure \(f \cdot \mu\) , where \(f\) is \(\mu\) -measurable, between 0 and 1. From this, deduce that \(\Lambda\) is \(\mu\) -adequate.
\end{enumerate}

\begin{enumerate}
\def\labelenumi{\arabic{enumi})}
\setcounter{enumi}{27}
\tightlist
\item
  Let \(\mathbf{X}\) be a locally compact space, \(\Lambda : t \mapsto \lambda_t\) a \(\mu\) -adequate mapping of \(\mathbf{T}\) into \(\mathcal{M}_+(\mathbf{X})\) . Let \(\mu' = g \cdot \mu\) be a positive measure with base \(\mu\) , such that \(\Lambda\) is scalarly essentially \(\mu'\) -integrable.
\end{enumerate}

\begin{enumerate}
\def\labelenumi{\alph{enumi})}
\item
  Show that \(\Lambda\) is \(\mu'\) -integrable (use Exer. 27 b), or the relation \(\mu' = \sup_n \left( \inf(\mu', n\mu) \right)\) and Exer. 9 of §3).
\item
  Show that \(\nu' = \int \lambda_t d\mu'(t)\) is a measure with base \(\nu = \int \lambda_t d\mu(t)\) .
\item
  Show that the mapping \(t \mapsto g(t)\lambda_t\) is \(\mu\) -adequate, and that its integral is \(\nu'\) .
\end{enumerate}

\begin{enumerate}
\def\labelenumi{\arabic{enumi})}
\setcounter{enumi}{28}
\tightlist
\item
  Let X be a locally compact space. Let \(\Lambda: t \mapsto \lambda_{t}\) be a \(\mu\) -adequate mapping of T into \(\mathcal{M}_{+}(X)\) . Let g be a function on X, positive, universally measurable, and locally bounded. Show that the family \(t \mapsto g \cdot \lambda_{t}\) is \(\mu\) -adequate, and that
\end{enumerate}

\[
\int (g \cdot \lambda_ {t}) d \mu (t) = g \cdot \int \lambda_ {t} d \mu (t)
\]

under each of the following conditions: a) \(g\) is moderated for the measure \(\int \lambda_t d\mu(t)\) ; b) the measures \(\lambda_t\) are bounded.

\begin{enumerate}
\def\labelenumi{\arabic{enumi})}
\setcounter{enumi}{29}
\tightlist
\item
  \begin{enumerate}
  \def\labelenumii{\alph{enumii})}
  \tightlist
  \item
    Let M be a bounded subset of C. Show that the closed convex envelope of M is the intersection of the closed disks containing M.
  \end{enumerate}
\end{enumerate}

\begin{enumerate}
\def\labelenumi{\alph{enumi})}
\setcounter{enumi}{1}
\tightlist
\item
  Let X be a locally compact space, \(\mu\) a complex measure on X such that \(\|\mu\| = \mu(1) = 1\) , f a bounded, continuous complex function on X. Show that \(\mu(f)\) belongs to the closed convex envelope of \(f(X)\) in C. (Use a.) Deduce from this a new proof of Prop. 9.
\end{enumerate}

\begin{enumerate}
\def\labelenumi{\arabic{enumi})}
\setcounter{enumi}{30}
\tightlist
\item
  Let \(\theta\) be a complex measure on \(T\) , and \(F\) a Banach space.
\end{enumerate}

\begin{enumerate}
\def\labelenumi{\alph{enumi})}
\tightlist
\item
  Show that the space \(\mathrm{L}_{\mathrm{loc}}^{1}(\mathrm{T},\theta ;\mathrm{F})\) is complete. (Write \(|\theta |\) in the form \(\sum_{\alpha}\mu_{\alpha}\) , where \((\mu_{\alpha})\) is a summable family of measures \(\geqslant 0\) on T whose supports form a locally countable family of pairwise disjoint compact subsets. This defines a continuous mapping of \(\mathrm{L}_{\mathrm{loc}}^{1}(\mathrm{T},\theta ;\mathrm{F})\) into \(\prod_{\alpha}\mathrm{L}_{\mathrm{F}}^{1}(\mathrm{T},\mu_{\alpha})\) . Let \(\mathfrak{F}\) be a Cauchy filter on \(\mathrm{L}_{\mathrm{loc}}^{1}(\mathrm{T},\theta ;\mathrm{F})\) . Its image
\end{enumerate}

in \(\prod_{\alpha} \mathrm{L}_{\mathrm{F}}^{1}(\mathrm{T}, \mu_{\alpha})\) converges to an element \((\mathbf{f}_{\alpha})\) . Show that the \(\mathbf{f}_{\alpha}\) define a \(\theta\) -measurable

function \(\mathbf{f}\) on \(\mathrm{T}\) , then that \(\mathbf{f}\) is locally \(\theta\) -integrable, and that \(\dot{\mathbf{f}}\) is the limit of \(\mathfrak{F}\) in \(\mathrm{L}_{\mathrm{loc}}^{1}(\mathrm{T},\theta ;\mathrm{F}).\)

\begin{enumerate}
\def\labelenumi{\alph{enumi})}
\setcounter{enumi}{1}
\tightlist
\item
  Let p be a real number \textgreater{} 1. Denote by \(\mathcal{L}_{\mathrm{loc}}^{p}(\mathrm{T},\theta;\mathrm{F})\) the vector space of \(\theta\) -measurable functions f on T with values in F, such that \(|f|^{p}\) is locally \(\theta\) -integrable. Equip it with the semi-norms \(\mathbf{f} \mapsto (\int |\mathbf{f}\varphi_{\mathrm{K}}|^{p}d|\theta|)^{1/p}\) , where K runs over the set of compact subsets of T. Let \(\mathrm{L}_{\mathrm{loc}}^{p}(\mathrm{T},\theta;\mathrm{F})\) be the associated Hausdorff space. Show that this space is compete. (Same method as in a).
\end{enumerate}

\setcounter{exercisegroup}{6}
\edef\CurrentExerciseGroup{\arabic{exercisegroup}}
\section*{\texorpdfstring{\S\CurrentExerciseGroup}{Section \CurrentExerciseGroup}}
\phantomsection
\addcontentsline{toc}{section}{\texorpdfstring{Exercises for \S\CurrentExerciseGroup}{Exercises for Section \CurrentExerciseGroup}}

\begin{enumerate}
\def\labelenumi{\arabic{enumi})}
\tightlist
\item
  Let \(\mathbf{T}\) and \(\mathbf{X}\) be two locally compact spaces, \(\mu\) a positive measure on \(\mathbf{T}\) , \(\pi\) a continuous and \(\mu\) -proper mapping of \(\mathbf{T}\) into \(\mathbf{X}\) , and let \(\nu = \pi(\mu)\) .
\end{enumerate}

\begin{enumerate}
\def\labelenumi{\alph{enumi})}
\item
  Show that for every numerical function \(g \geqslant 0\) , defined and lower semi-continuous on \(X\) , \(g \circ \pi\) is lower semi-continuous on \(T\) and \(\nu^{*}(g) = \mu^{*}(g \circ \pi)\) .
\item
  Show that if \(\mathbf{f}\) is a \(\nu\) -integrable function with values in a Banach space, then \(\mathbf{f} \circ \pi\) is \(\mu\) -integrable and \(\int \mathbf{f} d\nu = \int (\mathbf{f} \circ \pi) d\mu\) . Is the converse valid without a supplementary condition (cf.~§4, Exer. 2b)?
\end{enumerate}

¶ 2) Let T and X be two locally compact spaces, P the set of proper continuous mappings of T into X, equipped with the topology of compact convergence. Let H be a subset of P such that, for every compact subset K of X, the union of the sets \(\overline{\pi}^{-1}(K)\) , where \(\pi\) runs over H, is relatively compact in T.

\begin{enumerate}
\def\labelenumi{\alph{enumi})}
\item
  Show that the mapping \((\pi, \lambda) \mapsto \pi(\lambda)\) of \(\mathrm{H} \times \mathcal{M}(\mathrm{T})\) into \(\mathcal{M}(\mathrm{X})\) is not necessarily continuous when each of the spaces \(\mathcal{M}(\mathrm{T})\) and \(\mathcal{M}(\mathrm{X})\) is equipped with the quasi-strong topology (Ch. III, §1, Exer. 8).
\item
  If B is a bounded subset of \(\mathcal{M}(\mathrm{T})\) , show that the restriction to \(\mathrm{H} \times \mathrm{B}\) of the mapping \((\pi, \lambda) \mapsto \pi(\lambda)\) is continuous when each of the spaces \(\mathcal{M}(\mathrm{T}), \mathcal{M}(\mathrm{X})\) is equipped with the vague topology (make use of Prop. 15 of Ch. III, §1).
\item
  Show that the mapping \((\pi, \lambda) \mapsto \pi(\lambda)\) of \(\mathrm{H} \times \mathcal{M}_{+}(\mathrm{T})\) into \(\mathcal{M}_{+}(\mathrm{X})\) is continuous when \(\mathcal{M}(\mathrm{T})\) and \(\mathcal{M}(\mathrm{X})\) are equipped with the vague topology (cf.~Ch. III, §1, Exer. 10). d) Give an example where \(\mathrm{T} = \mathrm{X}\) is compact and the mapping \((\pi, \lambda) \mapsto \pi(\lambda)\) of \(\mathcal{P} \times \mathcal{M}(\mathrm{T})\) into \(\mathcal{M}(\mathrm{X}) = \mathcal{M}(\mathrm{T})\) is not continuous, when \(\mathcal{M}(\mathrm{T})\) is equipped with the vague topology (take for \(\mathrm{T}\) the torus \(\mathbf{T}\) and use Exer. 2 b) of Ch. III, §4).
\end{enumerate}

¶ 3) Let T be a noncompact locally compact space, \(\mu\) a positive measure on T, and \(\mathcal{G}\) the group of homeomorphisms of T. Give an example showing that the mapping \(\pi \mapsto \pi(\mu)\) of \(\mathcal{G}\) into \(\mathcal{M}(T)\) is not necessarily continuous when \(\mathcal{G}\) is equipped with the topology of compact convergence and \(\mathcal{M}(T)\) with the vague topology (take for T the subspace of R formed by 0 and the points n and 1/n (n an integer \(\geqslant 1\) ) and for \(\mu\) the measure defined by \(\mu(\{0\}) = 0\) , \(\mu(\{1/n\}) = 1/n^2\) , \(\mu(\{n\}) = 1\) ).

\begin{enumerate}
\def\labelenumi{\arabic{enumi})}
\setcounter{enumi}{3}
\item
  For every locally compact space T, let \(\mathcal{M}^{c}(T)\) be the subspace of \(\mathcal{M}(T)\) formed by the real measures with compact support. If T and X are two locally compact spaces, and \(\pi\) is a continuous mapping of T into X, then \(\pi\) is \(\lambda\) -proper and \(\pi(\lambda) \in \mathcal{M}^{c}(X)\) for every measure \(\lambda \in \mathcal{M}^{c}(T)\) . Give an example showing that the mapping \(\lambda \mapsto \pi(\lambda)\) of \(\mathcal{M}^{c}(T)\) into \(\mathcal{M}^{c}(X)\) is not necessarily continuous when \(\mathcal{M}^{c}(T)\) and \(\mathcal{M}^{c}(X)\) are equipped with the vague topology (or with the quasi-strong topology (Ch. III, §1, Exer. 8)) and that \(\pi\) is not a proper mapping (cf.~Ch. III, §2, Exer. 1).
\item
  \begin{enumerate}
  \def\labelenumii{\alph{enumii})}
  \tightlist
  \item
    Let \(\rho\) and \(\sigma\) be two positive measures on \(\mathbf{R}\) . Show that if \(\rho(]a,b]) = \sigma(]a,b])\) for every semi-open interval \(]a,b]\) , then \(\rho = \sigma\) .
  \end{enumerate}
\end{enumerate}

\begin{enumerate}
\def\labelenumi{\alph{enumi})}
\setcounter{enumi}{1}
\tightlist
\item
  Let I be an interval of \(\mathbf{R}\) , with left end-point \(\alpha\) and right end-point \(\beta\) , and let \(\psi\) be an increasing finite numerical function defined on I; extend \(\psi\) to \(\mathbf{R}\) in any way whatsoever. In order that \(\psi\) be proper for the measure \(\varphi_{\mathrm{I}} \cdot \mu\) ( \(\mu\) Lebesgue measure on \(\mathbf{R}\) ), it is necessary and sufficient that the following two conditions hold: \(1^{\circ}\) either \(\alpha\) is finite, or \(\alpha = -\infty\) and \(\lim_{x \to -\infty} \psi(x) = -\infty\) ; \(2^{\circ}\) either \(\beta\) is finite, or \(\beta = +\infty\) and
\end{enumerate}

\(\lim_{x\to +\infty}\psi (x) = +\infty .\)

Suppose these conditions verified. For every \(y \in \mathbf{R}\) , let \(\theta(y)\) be the supremum of the set of \(x \in I\) such that \(\psi(x) \leqslant y\) (the supremum being equal to \(\alpha\) if this set is empty). Show that \(\theta\) is a finite numerical function, increasing and right-continuous in \(\mathbf{R}\) . If \(\nu\) is the image of \(\varphi_{\mathrm{I}} \cdot \mu\) under \(\psi\) , then \(\nu(]a,b]) = \theta(b) - \theta(a)\) for every semi-open interval \(]a,b]\) in \(\mathbf{R}\) .

\begin{enumerate}
\def\labelenumi{\alph{enumi})}
\setcounter{enumi}{2}
\tightlist
\item
  Conversely, show that for every finite numerical function \(\theta\) , increasing and right-continuous in \(\mathbf{R}\) , there exists one and only one positive measure \(\nu\) on \(\mathbf{R}\) such that \(\nu(]a,b]) = \theta(b) - \theta(a)\) for every semi-open interval \(]a,b]\) of \(\mathbf{R}\) (the 'Stieltjes measure on \(\mathbf{R}\) defined by \(\theta'\) ; one writes \(\int f d\theta\) instead of \(\int f d\nu\) ); moreover, every positive measure on \(\mathbf{R}\) may be obtained in this way. Under what condition is the measure \(\nu\) diffuse? What is then the image of \(\nu\) under \(\theta\) ?
\end{enumerate}

\begin{enumerate}
\def\labelenumi{\arabic{enumi})}
\setcounter{enumi}{5}
\tightlist
\item
  \begin{enumerate}
  \def\labelenumii{\alph{enumii})}
  \tightlist
  \item
    Let K be the Cantor triadic set in the interval {[}0,1{]} (GT, IV, §2, No.~5). Show that there exist a diffuse positive measure \(\nu\) on \(\mathbf{R}\) , with support K and total mass 1, and a continuous increasing function \(\theta\) , defined on \(\mathbf{R}\) , such that \(\theta(\mathbf{R}) = \mathbf{I}\) and \(\theta(\nu) = \varphi_{\mathrm{I}} \cdot \mu\) ( \(\mu\) the Lebesgue measure; cf.~GT, IV, §8, Exer. 16).
  \end{enumerate}
\end{enumerate}

\begin{enumerate}
\def\labelenumi{\alph{enumi})}
\setcounter{enumi}{1}
\tightlist
\item
  Deduce from a) that there exists a diffuse positive measure on R, alien to Lebesgue measure, whose support is all of I (in each interval J contiguous to K and contained in I, take a measure proportional to the image of \(\nu\) under an affine mapping of I onto J; then proceed by recursion).
\end{enumerate}

\begin{enumerate}
\def\labelenumi{\arabic{enumi})}
\setcounter{enumi}{6}
\tightlist
\item
  Let \(\mu\) be Lebesgue measure on \(\mathbf{R}\) , I the interval \([0,1]\) of \(\mathbf{R}\) . If one sets \(\theta(x) = |x|\) , \(\theta\) is proper for the measure \(\varphi_{\mathrm{I}} \cdot \mu\) and one has \(\theta(\varphi_{\mathrm{I}} \cdot \mu) = \varphi_{\mathrm{I}} \cdot \mu\) . Give an example of a set negligible for \(\varphi_{\mathrm{I}} \cdot \mu\) , whose image under \(\theta\) is not measurable for \(\varphi_{\mathrm{I}} \cdot \mu\) (cf.~Ch. IV, §4, Exer. 8).
\end{enumerate}

¶ 8) Let T be a compact space, μ a diffuse positive measure on T of total mass 1.

\begin{enumerate}
\def\labelenumi{\alph{enumi})}
\tightlist
\item
  Show that there exists a continuous mapping \(\pi\) of T onto E = {[}0,1{]} such that the image of \(\mu\) under \(\pi\) is the Lebesgue measure \(\lambda\) on E. (Proceed as in the proof of Urysohn's theorem (GT, IX, §4, Th. 1) by defining, for every \(t\) such that \(0 \leqslant t \leqslant 1\) , an open set \(\mathrm{U}(t) \subset \mathrm{T}\) , quadrable for \(\mu\) (Ch. IV, §5, Exer. 17), such that \(\mathrm{U}(0) = \varnothing\) , \(\mathrm{U}(1) = \mathrm{T}\) , \(\overline{\mathrm{U}(t)} \subset \mathrm{U}(t')\) for \(t < t'\) , and finally \(\mu(\mathrm{U}(t)) = t\) ; one uses Exer. 7 of §5 to prove that if V, W are two quadrable open sets in T such that \(\overline{\mathrm{V}} \subset \mathrm{W}\) and \(\mu(\mathrm{V}) < \mu(\mathrm{W})\) , then there exists a quadrable open set U such that \(\overline{\mathrm{V}} \subset \mathrm{U} \subset \overline{\mathrm{U}} \subset \mathrm{W}\) and such that
\end{enumerate}

\[
\frac {1}{3} \mu (\mathrm{W-V}) \leqslant \mu (\mathrm{U-V}) \leqslant \frac {2}{3} \mu (\mathrm{W-V}).
\]

Finally, one applies Exer. 5 a).)

\begin{enumerate}
\def\labelenumi{\alph{enumi})}
\setcounter{enumi}{1}
\item
  Deduce from a) that there exist subsets of T that are not \(\mu\) -measurable (Ch. IV, §4, Exer. 8), that there exist sequences in \(\mathcal{L}^{p}(\mathrm{T},\mu)\) that convergence in mean of order p to 0 but do not converge to 0 at any point of T, for \(1 \leqslant p < +\infty\) (Ch. IV, §3, Exer. 1), and sequences \((f_{n})\) such that \((\widetilde{f}_{n})\) converges to 0 for the weakened topology of \(\mathrm{L}^{p}(\mathrm{T},\mu)\) but does not converge in measure to 0 (§5, Exer. 19).
\item
  Suppose, moreover, that T is metrizable. Show that there exist a \(\mu\) -negligible subset N of T, a \(\lambda\) -negligible subset M of E, and a homeomorphism \(\pi\) of E-M onto T-N such that, if one extends (arbitrarily) \(\pi\) to a mapping \(\varphi\) of E into T, and \(\pi^{-1}\) to a mapping \(\psi\) of T into E, then \(\varphi(\lambda)=\mu\) and \(\psi(\mu)=\lambda\) . (Using Exer. 17 of Ch. IV, §5, show that for every integer n\textgreater0, there exists a finite partition of T formed by a \(\mu\) -negligible set and by quadrable open sets, of diameter \(\leqslant1/n\) (for a metric compatible with the topology of T) and of measure \(\leqslant1/n\) . Proceeding recursively, deduce therefrom, by passage to the limit, the existence of a continuous mapping f of E-D into T, where D is a countable subset of E, such that \(f(\lambda)=\mu\) ; show that one can arrange that f be a homeomorphism of E-D onto a subset of T of measure 1 (cf.~§8, Exer. 14).)
\end{enumerate}

¶ 9) Let \(\mu\) and \(\nu\) be two diffuse positive measures on a locally compact space T. Show that, for \(\nu\) to be a measure with base \(\mu\) , it suffices that every \(\mu\) -measurable subset of T be also \(\nu\) -measurable (argue by contradiction using Th. 3 and Prop. 13 of §5, and Exer. 8 b) of §6).

\begin{enumerate}
\def\labelenumi{\arabic{enumi})}
\setcounter{enumi}{9}
\tightlist
\item
  In the interval \(\mathbf{I} = ]0,1[\) of \(\mathbf{R}\) , define the function \(g\) by \(g(t) = -\sqrt{n}\) for \(\frac{1}{2}\left(\frac{1}{n} + \frac{1}{n+1}\right) < t \leqslant \frac{1}{n}\) and \(g(t) = \sqrt{n}\) for \(\frac{1}{n+1} < t \leqslant \frac{1}{2}\left(\frac{1}{n} + \frac{1}{n+1}\right) (n = 1,2,\ldots)\) , and extend \(g\) by 0 in \(\mathbf{R} - \mathbf{I}\) ; \(g\) is integrable for Lebesgue measure. Set
\end{enumerate}

\[
\mathrm{G} (x) = \int_ {0} ^ {x} g (t) d t;
\]

show that the function \(f\) , equal to \(1 / \sqrt{x}\) on \([-1, +1]\) and to 0 elsewhere, which is integrable for Lebesgue measure, is such that \(t \mapsto f(\mathrm{G}(t))g(t)\) is not integrable on I.

¶ 11) The notations are those of No.~5; g is a locally μ-integrable numerical function on I.

\begin{enumerate}
\def\labelenumi{\alph{enumi})}
\item
  For G to be a \(\lambda\) -proper mapping of I into G(I), it is necessary and sufficient that the following conditions be fulfilled: \(1^{\circ}\) the limits G(a+) and G(b−) exist in \(\overline{\mathbf{R}}\) ; \(2^{\circ}\) if G(a+) ∈ G(I) then g is \(\mu\) -integrable on the interval {]} \(a, x_0[\) ; \(3^{\circ}\) if G(b−) ∈ G(I) then g is \(\mu\) -integrable on the interval {]} \(x_0, b[\) .
\item
  Suppose that the limits \(\mathrm{G}(a+)\) and \(\mathrm{G}(b-)\) exist in \(\overline{R}\) . Show that if f is such that \(t \mapsto \mathbf{f}(\mathrm{G}(t))g(t)\) is integrable for Lebesgue measure on I, then f is integrable on the interval with end-points \(\mathrm{G}(a+)\) and \(\mathrm{G}(b-)\) for Lebesgue measure, and formula (13) holds. (Reduce to the case that \(x_{0}\) is one of the numbers a, b; if, for example, \(x_{0} = a\) , observe that there exists in I a strictly increasing sequence \((b_{n})\) tending to b, such that the sequence \((\mathrm{G}(b_{n}))\) is either increasing or decreasing; apply Prop. 4 of Ch. IV, §4, No.~3 and Lebesgue's theorem).
\end{enumerate}

¶ 12) a) Let g be a finite numerical function defined on a compact interval I = {[}a, b{]} of R. One calls right-expansion set (resp. left-expansion set) of g in I the set of \(x \in I\) such that there exists \(y \in I\) for which x < y (resp. x \textgreater{} y) and \(g(x) < g(y)\) . Show that if g is continuous, then the right- (resp. left-) expansion set of g in I is open in I and that, if {]} \(\alpha, \beta[\) is a connected component \(^{1}\) of this set, then \(g(\alpha) = g(\beta)\) and \(g(x) \leqslant g(\alpha)\) for \(\alpha < x < \beta\) .

\begin{enumerate}
\def\labelenumi{\alph{enumi})}
\setcounter{enumi}{1}
\item
  Let \(r_1, r_2\) be two real numbers such that \(0 \leqslant r_1 < r_2\) . Let \(g\) be an increasing continuous function on I; let \(E'\) be the left-expansion set of \(g(x) - r_1x\) in I, \(E''\) the union of the right-expansion sets of \(g(x) - r_2x\) in each of the intervals that are the closures of the connected components of \(E'\) . If \(\mu\) is Lebesgue measure on I, show, using a), that \(\mu(E'') \leqslant \frac{r_1}{r_2} \mu(E')\) .
\item
  One calls the upper and lower right derivates at a point \(x \in I\) , of a finite numerical function \(f\) defined on \(I\) , the respective numbers
\end{enumerate}

\[
\begin{array}{l} \mathrm{D} _ {r} ^ {+} f (x) = \operatorname * {l i m s u p} _ {h \to 0,   h > 0} \left(f (x + h) - f (x)\right) / h \\ \mathrm{D} _ {r} ^ {-} f (x) = \operatorname * {l i m i n f} _ {h \to 0,   h > 0} \left(f (x + h) - f (x)\right) / h. \end{array}
\]

Similarly, one calls the upper and lower left derivates of \(f\) at the point \(x\) , the respective numbers

\[
\begin{array}{l}\mathrm{D} _ {l} ^ {+} f (x) = \underset {h \rightarrow 0, h <   0} {\lim \sup} \left(f (x + h) - f (x)\right) / h\\\mathrm{D} _ {l} ^ {-} f (x) = \underset {h \rightarrow 0, h <   0} {\lim \inf} \left(f (x + h) - f (x)\right) / h.\end{array}
\]

Show that if \(f\) is continuous and increasing, then the set of \(x \in I\) where \(\mathrm{D}_r^+ f(x) = +\infty\) and the set of \(x \in I\) where \(\mathrm{D}_l^- f(x) < \mathrm{D}_r^+ f(x)\) have measure zero for \(\mu\) (for every rational number \(r > 0\) (resp. for every pair \((r_1, r_2)\) of rational numbers such that \(0 \leqslant r_1 < r_2\) ), consider the set of \(x \in I\) such that \(\mathrm{D}_r^+ f(x) > r\) (resp. \(\mathrm{D}_l^- f(x) < r_1\) and \(\mathrm{D}_r^+ f(x) > r_2\) simultaneously), and apply \(b\) )). From this, deduce that for almost every \(x \in I\) , \(f\) admits a finite derivative (`Lebesgue's theorem') (apply the preceding result also to \(-f(-x)\) ).

¶ 13) In a compact interval I = {[}a, b{]} of R, let \((f_n)\) be a sequence of continuous, increasing finite functions, such that \(f_n(a) = 0\) for all n and such that the sum \(s(x) = \sum_{n=1}^{\infty} f_n(x)\) is finite on I. Show that there exists a set N ⊂ I, negligible for Lebesgue measure, such that, for every \(x \in I - N\) , the derivatives \(f_n'(x)\) and \(s'(x)\) exist and \(s'(x) = \sum_{n=1}^{\infty} f_n'(x)\) (`Fubini's theorem'). (Use Exer. 12. Observe that the series with general term \(f_n'(x)\) is convergent almost everywhere; setting

\[
s _ {n} (x) = \sum_ {k = 1} ^ {n} f _ {k} (x),
\]

consider next a subsequence \((s_{n_k})\) such that the series with general term \(s(x) - s_{n_k}(x)\) is convergent on I, and apply the preceding remark to this series.)

\begin{enumerate}
\def\labelenumi{\arabic{enumi})}
\setcounter{enumi}{13}
\item
  Let \(f\) be a numerical function defined on a compact interval \(I = [a, b]\) of \(\mathbf{R}\) , integrable on \(I\) for Lebesgue measure. If one sets \(F(x) = \int_{a}^{x} f(t) dt\) , show that \(F\) admits almost everywhere in \(I\) a derivative equal to \(f\) (reduce to the case that \(f \geqslant 0\) ; then, using Exer. 13, consider successively the case that \(f\) is lower semi-continuous then the general case (cf.~Ch. IV, §4, No.~4, Cor. of Th. 3)).
\item
  Denoting by \(\mu\) the Lebesgue measure on \(\mathbf{R}\) , a point \(x \in \mathbf{R}\) is said to be a density point of a subset A of \(\mathbf{R}\) if, when \(h\) and \(k\) tend to 0 through values \(>0\) , the quotient \(\mu^{*}(\mathrm{A} \cap [x - h, x + k]) / (h + k)\) tends to 1. Show that the set of points of an arbitrary subset A of \(\mathbf{R}\) that are not density points of A is negligible for \(\mu\) . (Reduce to the case that A is contained in a compact interval \([a, b]\) , and consider the function \(s_{\mathrm{A}}(x) = \mu^{*}(\mathrm{A} \cap [a, x])\) . Take a decreasing sequence \((\mathrm{A}_n)\) of open sets containing A, such that the series with general term \(s_{\mathrm{A}_n}(x) - s_{\mathrm{A}}(x)\) is convergent on \([a, b]\) , and use Exer. 13.)
\item
  \begin{enumerate}
  \def\labelenumii{\alph{enumii})}
  \tightlist
  \item
    Let \(g\) be a function integrable for Lebesgue measure on \(\mathbf{E} = [0,1]\) . Set \(\mathrm{G}(x) = \int_0^x g(t)dt\) ; show that at every point \(x \in \mathbf{E}\) where \(g\) is continuous, \(\mathbf{G}\) admits a derivative equal to \(g\) .
  \end{enumerate}
\end{enumerate}

\begin{enumerate}
\def\labelenumi{\alph{enumi})}
\setcounter{enumi}{1}
\item
  Give an example of a function \(g\) that is bounded and is continuous almost everywhere in \(\mathbf{E}\) , such that \(\mathbf{G}\) has no right derivative at the points of an uncountable subset of \(\mathbf{E}\) (cf.~FRV, I, §2, Exer. 8).
\item
  Show that the function equal to \(x^2 \sin \frac{1}{x^2}\) for \(x \neq 0\) , and to 0 for \(x = 0\) , admits a derivative at every point of \(E\) , but this derivative is not integrable for Lebesgue measure on \(E\) .
\end{enumerate}

\begin{enumerate}
\def\labelenumi{\arabic{enumi})}
\setcounter{enumi}{16}
\item
  Let T and X be two locally compact spaces, \(\pi\) a continuous mapping of T into X. Show that if \(\pi\) is \(\mu\) -proper for every positive measure \(\mu\) on T, then \(\pi\) is proper. (Argue by contradiction, using Prop. 7 of GT, I, §10.)
\item
  In Prop. 4 b), the conclusion may fail if one omits the hypothesis that \(\pi'\) is continuous. (Take \(T = R\) , \(T' = \overline{R}\) , \(T'' = R\) ; take for \(\mu\) Lebesgue measure, for \(\pi\) the canonical injection; set \(\pi'(x) = x\) for \(x \in R\) and \(\pi'(\pm \infty) = 0\) .)
\item
  Let \(\mathrm{T}, \mathrm{X}, \mathrm{Y}\) be locally compact spaces, \(\pi : \mathrm{T} \to \mathrm{X}\) , \(\pi' : \mathrm{X} \to \mathrm{Y}\) mappings, \(\pi'' = \pi' \circ \pi\) , \(\mu\) a real measure on \(\mathrm{T}\) . It can happen that \(\pi\) is \(\mu\) -proper and \(\pi'\) is \(\pi(\mu)\) -proper, without \(\pi''\) being \(\mu\) -proper. (Take \(\mathrm{T} = \mathbf{R} \times \{0, 1\}\) , \(\mathrm{X} = \mathbf{R}\) , \(\mathrm{Y} = \{0\}\) , \(\pi = \mathrm{pr}_1\) ; take for \(\mu\) the measure that induces Lebesgue measure (resp. the opposite of Lebesgue measure) on \(\mathbf{R} \times \{0\}\) (resp. \(\mathbf{R} \times \{1\}\) ) canonically identified with \(\mathbf{R}\) .)
\item
  Let \(\mathrm{T},\mathrm{X},\mathrm{Y}\) be three locally compact spaces, \(\mu\) a positive measure on \(\mathrm{T}\) , \(t\mapsto \lambda_t\) a \(\mu\) -adequate mapping of \(\mathrm{T}\) into \(\mathcal{M}_{+}(\mathrm{X})\) , \(\nu = \int \lambda_t d\mu (t)\) , and \(\pi\) a \(\nu\) -proper mapping of \(\mathrm{X}\) into \(\mathrm{Y}\) . Assume one of the following hypotheses: a) \(\mathrm{Y}\) is countable at infinity and \(\nu\) is moderated; b) \(\mathrm{Y}\) is countable at infinity and \(\lambda_t\) is bounded locally almost everywhere in \(\mathrm{T}\) . Then \(\pi\) is \(\lambda_t\) -proper locally almost everywhere in \(\mathrm{T}\) , the mapping \(t\mapsto \pi (\lambda_t)\) is \(\mu\) -adequate, and its integral is equal to \(\pi (\mu)\) .
\item
  Let \(\mathrm{T},\mathrm{X},\mathrm{Y}\) be three locally compact spaces, \(\pi : \mathrm{T} \to \mathrm{X}\) , \(\pi': \mathrm{X} \to \mathrm{Y}\) universally measurable mappings. Then \(\pi' \circ \pi\) is universally measurable. (Let \(\mu\) be a positive measure on \(\mathrm{T}\) . To prove that \(\pi' \circ \pi\) is \(\mu\) -measurable, reduce to the case that \(\mathrm{T}\) is compact and \(\pi\) is continuous, and use the fact that \(\pi'\) is measurable for \(\pi(\mu)\) .)
\item
  Let \(X, Y\) be two compact spaces, \(U\) a continuous linear mapping of the Banach space \(\mathcal{C}(X; \mathbf{R})\) into the Banach space \(\mathcal{C}(Y; \mathbf{R})\) , \(\pi : X \to Y\) a continuous mapping. Assume that for every \(y \in Y\) and every function \(f \in \mathcal{C}(X; \mathbf{R})\) ,
\end{enumerate}

\[
\inf _ {\pi (x) = y} f (x) \leqslant (U \cdot f) (y) \leqslant \sup _ {\pi (x) = y} f (x).
\]

Show that under these conditions there exists a vaguely continuous mapping \(\sigma: y \mapsto \sigma_y\) of Y into \(\mathcal{M}_+(X)\) such that, for every \(y \in Y\) , \(\sigma_y\) is carried by \(\overline{\pi}^{-1}(y)\) and one has \((U \cdot f)(\pi(x)) = \langle f, \sigma_{\pi(x)} \rangle\) for all \(x \in X\) ; converse. (Consider the transpose \(^t U\) .)

\setcounter{exercisegroup}{7}
\edef\CurrentExerciseGroup{\arabic{exercisegroup}}
\section*{\texorpdfstring{\S\CurrentExerciseGroup}{Section \CurrentExerciseGroup}}
\phantomsection
\addcontentsline{toc}{section}{\texorpdfstring{Exercises for \S\CurrentExerciseGroup}{Exercises for Section \CurrentExerciseGroup}}

\begin{enumerate}
\def\labelenumi{\arabic{enumi})}
\item
  Let X be a locally compact subspace of a locally compact space T. Show that for every measure \(\lambda\) on T, the support of the measure \(\lambda_{X}\) is the trace on X of the support of the measure \(\varphi_{X} \cdot \lambda\) .
\item
  Let \(\mathrm{T}\) be a locally compact space, \(\mathrm{X}\) a locally compact subspace of \(\mathrm{T}\) , and \(\mu\) a positive measure on \(\mathrm{T}\) ; set \(\nu = \varphi_{\mathrm{X}} \cdot \mu = j(\mu_{\mathrm{X}})\) ( \(j\) the canonical injection of \(\mathrm{X}\) into \(\mathrm{T}\) ).
\end{enumerate}

\begin{enumerate}
\def\labelenumi{\alph{enumi})}
\item
  Show that every \(\mu_{\mathrm{X}}\) -negligible subset of \(\mathbf{X}\) is \(\nu\) -negligible.
\item
  For every numerical function \(g \geqslant 0\) defined on \(X\) , show that \(\int^{*} g \, d\mu_{X} = \int_{X}^{*} g \, d\nu\) (use \(a\) ), as well as formula (7) of §3).
\item
  Let \(\mathbf{g}\) be a function defined on \(X\) , with values in \(\overline{\mathbf{R}}\) or in a Banach space, \(\mathbf{g}'\) the extension of \(\mathbf{g}\) to \(T\) equal to 0 on \(T - X\) . For \(\mathbf{g}\) to be \(\mu_{X}\) -integrable, it is necessary and sufficient that \(\mathbf{g}'\) be \(\nu\) -integrable, in which case \(\int \mathbf{g} d\mu_X = \int_X \mathbf{g} d\nu\) (use \(b\) )).
\end{enumerate}

\begin{enumerate}
\def\labelenumi{\arabic{enumi})}
\setcounter{enumi}{2}
\tightlist
\item
  The notations being the same as in Exer. 2, assume that \(\mathbf{X}\) is an open subset of \(\mathbf{T}\) .
\end{enumerate}

\begin{enumerate}
\def\labelenumi{\alph{enumi})}
\item
  For every numerical function \(g \geqslant 0\) defined on \(X\) , show that \(\int^{*} g d\mu_{X} = \int_{X}^{*} g d\mu\) (use Exer. 2, Prop. 3 of §5 and Prop. 4 of §1).
\item
  Let g be a function defined on X, with values in \(\overline{R}\) or in a Banach space, \(g'\) the extension of g to T equal to 0 on T-X. For g to be \(\mu_{X}\) -integrable, it is necessary and sufficient that \(g'\) be \(\mu\) -integrable, in which case \(\int g d\mu_{X} = \int_{X} g d\mu\) (use a)).
\end{enumerate}

\begin{enumerate}
\def\labelenumi{\arabic{enumi})}
\setcounter{enumi}{3}
\item
  The notations being those of Exer. 2, show that if X is closed in T then, for every numerical function \(f \geqslant 0\) defined on T, \(\int^{*}(f \circ j) d\mu_{X} = \int^{*} f d\nu\) (observe that the mapping \(j\) is proper, and use Prop. 2 of §4). Under the same conditions, if \(\mathbf{f}\) is a mapping of T into \(\overline{\mathbf{R}}\) or into a Banach space, for \(\mathbf{f} \circ j\) to be \(\mu_{X}\) -integrable it is necessary and sufficient that \(\mathbf{f}\) be \(\nu\) -integrable, in which case \(\int(\mathbf{f} \circ j)\mu_{X} = \int \mathbf{f} d\nu\) (§4, Th. 2).
\item
  Let \(\mathbf{T}\) and \(\mu\) be the locally compact space and the measure defined in Exer. 5 of Ch. IV, §1, and let \(\mathbf{D}\) be the (closed) set of points of \(\mathbf{T}\) of the form \((0,y)\) .
\end{enumerate}

\begin{enumerate}
\def\labelenumi{\alph{enumi})}
\item
  Show that the measure \(\mu_{\mathrm{D}}\) induced on D by \(\mu\) is zero; however, \(\int^{*}\varphi_{\mathrm{D}}d\mu = +\infty\) (cf.~Exer. 3 a)).
\item
  Let \(\mathrm{X} = \mathrm{T} - \mathrm{D}\) , and let \(j\) be the canonical injection of \(\mathrm{X}\) into \(\mathrm{T}\) ; then \(j(\mu_{\mathrm{X}}) = \mu\) , but \(\int^{*} (\varphi_{\mathrm{D}} \circ j) d\mu_{\mathrm{X}} = 0\) , \(\int^{*} \varphi_{\mathrm{D}} d\mu = +\infty\) (cf.~Exer. 4).
\end{enumerate}

\begin{enumerate}
\def\labelenumi{\arabic{enumi})}
\setcounter{enumi}{5}
\item
  Let T be a locally compact space, X a locally compact subspace of T, and j the canonical injection of X into T. Let \(\lambda\) be a positive measure on X such that, for every compact set \(K \subset T\) , \(K \cap X\) is essentially \(\lambda\) -integrable. Show that the measure \(\nu = j(\lambda)\) is the smallest positive measure \(\rho\) on T such that \(\rho_{X} = \lambda\) , and that its support is the closure in T of the support of \(\lambda\) .
\item
  Let X, Y be two locally compact spaces, \(\pi: X \to Y\) a continuous mapping, T a locally compact subspace of Y, \(S = \overline{\pi}^{-1}(T)\) , \(\pi_{T}: S \to T\) the mapping that coincides with \(\pi\) in S, and \(\nu\) a positive measure on Y. Show that the following three conditions are equivalent:
\end{enumerate}

\begin{enumerate}
\def\labelenumi{\alph{enumi})}
\item
  If \(\nu_{\mathrm{T}}\) is the measure induced by \(\nu\) on \(\mathbf{T}\) , there exists a positive measure \(\mu_{\mathrm{T}}\) on \(\mathbf{S}\) such that \(\pi_{\mathrm{T}}(\mu_{\mathrm{T}}) = \nu_{\mathrm{T}}\) .
\item
  There exists a positive measure \(\lambda\) on \(X\) such that \(\pi (\lambda) = \nu_{\mathrm{T}}\) .
\item
  There exists a positive measure \(\lambda_{\mathrm{S}}\) on S such that \((\pi|\mathrm{S})(\lambda_{\mathrm{S}})=\nu_{\mathrm{T}}\) .
\end{enumerate}

\begin{enumerate}
\def\labelenumi{\arabic{enumi})}
\setcounter{enumi}{7}
\item
  Let X, Y be two locally compact spaces, \(\pi: X \to Y\) a continuous mapping, and \(\nu\) a positive measure on Y that is concentrated on \(\pi(X)\) . Assume that for every \(y \in \pi(X)\) , there exist a neighborhood V of y in Y and a positive measure \(\lambda_{\mathrm{V}}\) on \(\overline{\pi}^{-1}(V)\) such that the image of \(\lambda_{V}\) under the mapping \(\pi_{\mathrm{V}}: \overline{\pi}^{-1}(V) \to V\) that coincides with \(\pi\) is equal to the induced measure \(\nu_{V}\) . Show that under these conditions, there exists a positive measure \(\mu\) on X such that \(\pi(\mu) = \nu\) . (Consider the set G of pairs \((G, \lambda)\) , where G is open in Y and \(\lambda\) is a positive measure on X such that \(\pi(\lambda) = \nu_{\mathrm{G}}\) . Order G by the relation \(\langle G_{1} \subset G_{2} \text{ and } \lambda_{1} \leqslant \lambda_{2} \rangle\) between \((G_{1}, \lambda_{1})\) and \((G_{2}, \lambda_{2})\) . First show that the ordered set G is inductive, then complete the argument with the help of Zorn's lemma and Exer. 7.)
\item
  Let X, Y be two locally compact spaces, \(\pi: X \to Y\) a continuous mapping. The mapping \(\pi\) is said to be conservative if, for every positive measure \(\nu\) on Y that is concentrated on \(\pi(X)\) , there exists a positive measure \(\mu\) on X such that \(\nu = \pi(\mu)\) .
\end{enumerate}

\begin{enumerate}
\def\labelenumi{\alph{enumi})}
\item
  Let X,Y,Z be three locally compact spaces, \(\pi:X\to Y\) and \(\pi^{\prime}:Y\to Z\) two continuous mappings; set \(\pi^{\prime\prime}=\pi^{\prime}\circ\pi\) and assume that the mapping \(\pi\) is surjective. Show that if \(\pi\) and \(\pi^{\prime}\) are conservative, then so is \(\pi^{\prime\prime}\) ; conversely, if \(\pi^{\prime\prime}\) is conservative then so is \(\pi^{\prime}\) .
\item
  Assume that for every \(y \in \pi(X)\) there exists an open neighborhood V of y in Y such that, if \(\pi_{\mathrm{V}} : \overline{\pi}^{-1}(\mathrm{V}) \to \mathrm{V}\) is the mapping that coincides with \(\pi\) , there exists a continuous section associated with \(\pi_{V}\) (GT, I, §3, No.~5). Show that \(\pi\) is conservative (make use of Exer. 8).
\item
  Let Y be the interval \([0,1]\) of R, X the set Y equipped with the discrete topology, and \(\pi:X\to Y\) the identity mapping. Show that the mapping \(\pi\) is not conservative.
\end{enumerate}

\begin{enumerate}
\def\labelenumi{\arabic{enumi})}
\setcounter{enumi}{9}
\item
  Let \(X, Y\) be two locally compact spaces. Show that every proper continuous mapping \(\pi\) of \(X\) into \(Y\) is conservative. (The hypothesis implies that the mapping \(f \mapsto f \circ \pi\) of \(\mathcal{K}(Y)\) into \(\mathcal{C}(X)\) has its image \(E\) contained in \(\mathcal{K}(X)\) ; then, for every positive measure \(\nu\) on \(Y\) , apply to the subspace \(E\) of \(\mathcal{K}(X)\) and to the positive linear form \(f \circ \pi \mapsto \nu(f)\) on \(E\) the Cor. of Prop. 1 of TVS, II, §3, No.~1.)
\item
  Let X, Y be two locally compact spaces, \(\pi: X \to Y\) a continuous mapping, and M a subset of Y such that \(\overline{\pi}^{-1}(M)\) is contained in a countable union of compact sets. Then, for every positive measure \(\nu\) on Y that is concentrated on M, there exists a positive measure \(\mu\) on X such that \(\pi(\mu) = \nu\) (use Exer. 10).
\end{enumerate}
